\documentclass[11pt]{article}

\usepackage[margin=1in]{geometry}
\usepackage{amsmath,amssymb}
\usepackage[hidelinks]{hyperref}

\title{Literature Answer to Question 3 of arXiv:1509.01499}
\author{fa\_banach\_001, agent\_lane\_12}
\date{2026-06-20}

\begin{document}
\maketitle

\section*{Status}

This is a lightweight \texttt{literature\_already\_answered} packet. It records
that Question 3 in the final open-question section of Flores, Hernandez, and
Tradacete, \emph{Disjointly homogeneous Banach lattices and applications},
arXiv:1509.01499, was answered by a separate later paper.

\section*{Source Question}

The source paper asks, in Section ``Some open questions'':
\[
\text{Is there a reflexive }p\text{-DH r.i. space on }[0,1]\text{ whose dual is not DH?}
\]
Here \(p\)-DH means \(p\)-disjointly homogeneous and r.i. means
rearrangement invariant.

\section*{Supporting Answer}

The supporting paper is Sergey V. Astashkin,
\emph{Duality problem for disjointly homogeneous rearrangement invariant
spaces}, arXiv:1805.00691.

Its introduction explicitly identifies the same problem, saying that it solves
the duality problem and referring to Question 3 in the survey. The decisive
result is Theorem 2:
\[
\text{For every }1<p<\infty\text{ there exists a }p\text{-DH reflexive r.i.
space }X_p\text{ on }[0,1]\text{ such that }X_p^*\text{ is not DH.}
\]
This gives a positive answer to the source question.

\section*{Scope}

Only Question 3 from the source paper's final list is recorded as answered
here. The source paper's other final questions, including the complemented
disjoint sequence question, the separable DH versus DC question, the symmetric
atomic lattice question, and the Kato-property question, are outside this
packet.

\section*{References}

\begin{enumerate}
\item J. Flores, F. L. Hernandez, and P. Tradacete,
\emph{Disjointly homogeneous Banach lattices and applications},
arXiv:1509.01499, 2015.
\item S. V. Astashkin,
\emph{Duality problem for disjointly homogeneous rearrangement invariant
spaces}, arXiv:1805.00691, 2018.
\end{enumerate}

\end{document}
